% introduction.tex

% Section Title
\section{INTRODUCTION} \label{sec:introduction}

    Dynamic malware analysis leverages runtime behavior to detect malicious software, circumventing the limitations of static signature-based methods against obfuscated or polymorphic variants. API call sequences---the ordered invocations a program makes to the operating system---provide a rich behavioral fingerprint, capturing memory operations, registry accesses, and system queries that reveal malicious intent.

    This laboratory implements a comparative study of neural network architectures for binary malware classification based on API call sequences. The dataset consists of execution traces, each represented as a variable-length sequence of API function names (e.g., \cooltext{NtAllocateVirtualMemory}, \cooltext{RegOpenKeyExW}).

    We evaluate four progressively complex approaches:

    \begin{itemize}
        \item \textbf{1:} Implement and evaluate a frequency-based logistic regression baseline.
        \item \textbf{2:} Develop Feed-Forward Neural Networks with different input representations.
        \item \textbf{3:} Explore Recurrent Neural Network architectures for sequential modeling.
        \item \textbf{4:} Construct Graph Neural Networks to leverage structural information in API call sequences.
    \end{itemize}

    The following sections detail the methodology and results for each task, analyzing how architectural choices impact classification performance.
